\bookchapter{الفصل 46: طريق التقدم — من مستخدم إلى مهندس بنية}

\chapterintro{حين تتقن منهجيات الفصول 37 إلى 44 وتصبح قادرا على إنجاز تطوير المشاريع بالذكاء الاصطناعي بمهارة، فما الخطوة التالية؟ يستكشف هذا الفصل طريق التقدم من “مستخدم لأدوات الذكاء الاصطناعي” إلى “مهندس بنية لأنظمة الذكاء الاصطناعي”. لم يعد التركيز في هذه المرحلة على “كيف أجعل الذكاء الاصطناعي يكتب كودا جيدا”، بل على “كيف أدمج الذكاء الاصطناعي بسلاسة أكبر في دورة التطوير بأكملها”.}

\booksection{أولاً: بروتوكول MCP وتوسيع الأدوات}

بروتوكول MCP (Model Context Protocol) بروتوكول مفتوح ظهر عام 2025 وانتشر تدريجيا في 2026، ويحدد معيار الاتصال بين نماذج الذكاء الاصطناعي والأدوات الخارجية. بعبارة بسيطة، يتيح MCP للذكاء الاصطناعي “استخدام الأدوات” — فلا يقتصر على توليد الكود، بل ينفذ الأوامر ويستعلم قواعد البيانات ويستدعي واجهات API ويتعامل مع نظام الملفات.

قبل MCP كان لكل أداة ذكاء اصطناعي أسلوبها الخاص في دمج الأدوات، ولا توافق بينها. ومع ظهور MCP أصبح ممكنا أن تستخدم أي أداة مطابقة لمعيار MCP في أي أداة ذكاء اصطناعي تدعمه. هذا خفض كثيرا تكلفة دمج الأدوات، ووسع حدود قدرات الذكاء الاصطناعي من “توليد النصوص” إلى “تنفيذ الإجراءات”.

بالنسبة للمستخدم المتقدم، تكمن قيمة MCP في إمكانية بناء قدرات أدوات مخصصة للذكاء الاصطناعي. مثلا إذا كان فريقك يملك نظام نشر داخليا، يمكنك كتابة أداة MCP تتيح للذكاء الاصطناعي استدعاء نظام النشر مباشرة دون الانتقال اليدوي إلى الطرفية. وهذه قدرة “استدعاء الذكاء الاصطناعي للأدوات” هي البنية التحتية للهندسة الوكيلة (Agentic Engineering).

\booksection{ثانياً: سير العمل المخصص ومهارات Skills}

حين تكرر تنفيذ بعض مسارات البرمجة بالذكاء الاصطناعي الثابتة، يمكنك تثبيتها في شكل “سير عمل مخصص” أو “Skills”. مثلا قد يملك فريقك مسار مراجعة كود معياريا — فحص الأمان أولا ثم الأداء ثم قابلية الصيانة. يمكنك كتابة هذا المسار في شكل Skill لينفذه الذكاء الاصطناعي تلقائيا في كل مراجعة كود.

القيمة الجوهرية لسير العمل المخصص هي “الاتساق” — فهو يضمن اتباع المعيار نفسه في كل تنفيذ، دون انحراف بسبب اختلاف طول الحوار أو السياق. وهذا مهم بوجه خاص للعمل الجماعي: حين يستخدم كل عضو سير العمل نفسه، تتقارب أسلوب الكود الناتج ومعايير جودته.

تختلف طرق بناء سير العمل المخصص باختلاف الأداة: يدعم Cursor الـ Rules والتعليمات المخصصة (Custom Instructions)، ويدعم Claude Code الأوامر المخصصة، ويدعم Replit الـ Runme. لكن الفكرة الجوهرية واحدة: “ترميز” أفضل ممارساتك في قوالب تعليمات قابلة للتنفيذ من الذكاء الاصطناعي.

\booksection{ثالثاً: نمط التعاون متعدد الوكلاء}

بحلول 2026 أصبح التعاون متعدد الوكلاء موضوعا ساخنا. فكرته الجوهرية: لا نكتفي بذكاء اصطناعي واحد لإنجاز المهمة، بل نستخدم عدة وكلاء بأدوار مختلفة يتعاونون فيما بينهم. مثلا وكيل يكتب الكود، ووكيل يكتب الاختبارات، ووكيل يراجع الكود — لكل منهم دوره ويتحققون من عمل بعضهم.

تكمن ميزة التعاون متعدد الوكلاء في “التقسيم” و“التحقق المتقاطع”: كل وكيل يتخصص في مجاله فترتفع جودة المخرجات، والتحقق المتبادل بين الوكلاء يكشف المشاكل التي قد تغفل عنها وحدك. لكن العيوب واضحة أيضا: تكلفة التنسيق مرتفعة، والنفقات التواصلية كبيرة، والتصحيح أعقد.

\begin{tcbinfo}{نصيحة: يمكنك ممارسة "الوكلاء المتعددين" وحدك}
بصفتك مطورا فرديا لا تحتاج إلى بناء نظام وكلاء متعددين معقد — يمكنك “أداء” أدوار مختلفة لمحاكاة التعاون متعدد الوكلاء. مثلا: افتح حوارا ليكتب الذكاء الاصطناعي كود الوظائف؛ ثم افتح حوارا جديدا (مع توفير سياق CLAUDE.md) واطلب منه “مراجعة هذا الكود بصفته مراجع أكواد”؛ وأخيرا افتح حوارا ثالثا واطلب منه “تصحيح الكود بناء على ملاحظات المراجعة”. بهذا أسلوب “شخص واحد بأدوار متعددة” تحصل على معظم فوائد التعاون متعدد الوكلاء دون بنية تحتية معقدة.

\end{tcbinfo}

\booksection{رابعاً: مسار التعلم المستمر}

البرمجة بالذكاء الاصطناعي مجال سريع التغير. ما يعد أفضل ممارسة اليوم قد يحتاج إلى تحديث غدا. إليك بعض النصائح للتعلم المستمر:

أولا، تابع المصادر الأصلية لا القراءات الثانوية. تغريدات Andrej Karpathy الأصلية وملف CLAUDE.md، والوثائق الرسمية لـ Anthropic وOpenAI، وسجلات تحديثات الأدوات الكبرى — هذه هي المعلومات من مصدرها الأول. أما مقالات المجتمع والدروس فهي قراءات وملخصات لهذه المعلومات، لها قيمة لكنها قد تتشوه. اكتسب عادة قراءة المصادر الأصلية.

ثانيا، تعلم مدفوعا بالممارسة. لا “تتعلم ثم تنجز” بل “أنجز وأنت تتعلم”. كلما واجهت نوع مشكلة جديدا (مثل أول مرة تعالج فيها رفع ملفات، أو أول اتصال بواجهة API خارجية)، خذه فرصة لتعلم المعرفة ذات الصلة. كفاءة التعلم في الممارسة تفوق بكثير التعلم النظري المنفصل عن السياق.

ثالثا، حافظ على حساسيتك للأدوات دون مطاردة مفرطة للجديد. كل شهر أو شهرين خذ نصف يوم للتعرف على أحدث الأدوات والاتجاهات — لكن لا تتحول إلى كل أداة جديدة تظهر. معيار الحكم: هل تعالج الأداة الجديدة فعلا أحد اختناقاتك بشكل ملحوظ؟ إن لم تكن كذلك فالمزيج الحالي كاف.

رابعا، ابن قاعدة معرفتك الشخصية. سجل كل حوار حللت فيه مشكلة معقدة، وكل درس تعلمته من الوقوع في الأخطاء، وكل طريقة أثبتت فعاليتها. لا حاجة إلى النشر العلني — هذه “مذكرة منهجيتك” الخاصة. ومع تراكم الوقت ستصبح أثمن أصولك.

\begin{tcbpractice}{لنطبق عملياً: افتح "دفتر أخطاء وتجارب"}
أنشئ اليوم وثيقة (برنامج ملاحظات أو ملف .md كلاهما يصلح)، واكتب فيها أهم ثلاث خلاصات في البرمجة بالذكاء الاصطناعي تود حفظها حتى الآن — حتى لو كانت “تذكر أن تفكر جيدا قبل أن تسأل”. وبعد كل خطأ تقع فيه أو حيلة جديدة تتعلمها، أضف ملاحظة جديدة. وبعد نصف عام، حين تعود لتقرأها، ستشكر نفسك على ما فعلته اليوم.

\end{tcbpractice}

\begin{tcbsummary}{خلاصة الفصل}
\begin{itemize}
  \item بروتوكول MCP هو معيار الاتصال بين الذكاء الاصطناعي والأدوات: يتيح للذكاء الاصطناعي تنفيذ الأوامر واستعلام قواعد البيانات واستدعاء واجهات API
  \item القيمة الجوهرية لسير العمل المخصص و Skills هي "الاتساق": ترميز أفضل الممارسات في قوالب تعليمات
  \item الوكلاء المتعددون = تقسيم أدوار + تحقق متقاطع؛ ويمكن للمطور الفردي محاكاته بأسلوب "شخص واحد بأدوار متعددة"
  \item حيل التعلم المستمر الأربع: قراءة المصادر الأصلية، التعلم أثناء الممارسة، عدم المطاردة المفرطة للجديد، بناء قاعدة معرفة شخصية
\end{itemize}

\end{tcbsummary}
